
\documentclass{article}
\usepackage{colortbl}
\usepackage{makecell}
\usepackage{multirow}
\usepackage{supertabular}

\begin{document}

\newcounter{utterance}

\centering \large Interaction Transcript for game `cladder', experiment `full\_v1.5\_default', episode 2060 with selfrefine{-}Qwen3.5{-}27B.
\vspace{24pt}

{ \footnotesize  \setcounter{utterance}{1}
\setlength{\tabcolsep}{0pt}
\begin{supertabular}{c@{$\;$}|p{.15\linewidth}@{}p{.15\linewidth}p{.15\linewidth}p{.15\linewidth}p{.15\linewidth}p{.15\linewidth}}
    \# & \multicolumn{2}{c}{Player} && \multicolumn{2}{c}{Game Master} \\
    \hline

    \theutterance \stepcounter{utterance}  
    & & & \multicolumn{4}{p{0.6\linewidth}}{
        \cellcolor[rgb]{0.9,0.9,0.9}{
            \makecell[{{p{\linewidth}}}]{
                \texttt{\tiny{[P1$\langle$GM]}}
                \texttt{You are an expert in causal inference. The following question is not a typical commonsense query, but rather a meticulously designed question created by a professor specializing in causal inference, intended to assess the students' mastery of the course content.} \\
\texttt{Imagine a self{-}contained, hypothetical world with only the following conditions, and without any unmentioned factors or causal relationships: Encouragement level has a direct effect on studying habit and exam score. Studying habit has a direct effect on exam score. The overall probability of encouragement is 91\%. For students who are not encouraged, the probability of high exam score is 72\%. For students who are encouraged, the probability of high exam score is 52\%. Is high exam score more likely than low exam score overall?} \\
\texttt{Start your answer with 'Yes' or 'No', followed by additional reasoning or evidence to support your explanation.} \\
            }
        }
    }
     \\ \\

    \theutterance \stepcounter{utterance}  
    & \multicolumn{4}{p{0.6\linewidth}}{
        \cellcolor[rgb]{0.561,0.737,0.561}{
            \makecell[{{p{\linewidth}}}]{
                \texttt{\tiny{[P1$\rangle$GM]}}
                \texttt{Yes} \\
\texttt{To determine whether a high exam score is more likely than a low exam score overall, we calculate the total probability of a high exam score, \$P(\textbackslash{}text\{High Score\})\$, using the Law of Total Probability.} \\
\\ 
\texttt{Let \$E\$ represent the event of being encouraged, and \$S\$ represent the event of a high exam score.} \\
\texttt{From the problem description, we have:} \\
\texttt{1.  \$P(E) = 0.91\$} \\
\texttt{2.  \$P(\textbackslash{}neg E) = 1 {-} 0.91 = 0.09\$} \\
\texttt{3.  \$P(S | \textbackslash{}neg E) = 0.72\$} \\
\texttt{4.  \$P(S | E) = 0.52\$} \\
\\ 
\texttt{Using the Law of Total Probability:} \\
\texttt{\$\$P(S) = P(S | E) \textbackslash{}cdot P(E) + P(S | \textbackslash{}neg E) \textbackslash{}cdot P(\textbackslash{}neg E)\$\$} \\
\texttt{\$\$P(S) = (0.52 \textbackslash{}times 0.91) + (0.72 \textbackslash{}times 0.09)\$\$} \\
\texttt{\$\$P(S) = 0.4732 + 0.0648\$\$} \\
\texttt{\$\$P(S) = 0.538\$\$} \\
\\ 
\texttt{The overall probability of a high exam score is \$0.538\$ (or \$53.8\textbackslash{}\%\$).} \\
\texttt{The probability of a low exam score is \$1 {-} 0.538 = 0.462\$ (or \$46.2\textbackslash{}\%\$).} \\
\\ 
\texttt{Since \$0.538 > 0.462\$, a high exam score is more likely than a low exam score overall.} \\
            }
        }
    }
    & & \\ \\

    \theutterance \stepcounter{utterance}  
    & & & \multicolumn{2}{p{0.3\linewidth}}{
        \cellcolor[rgb]{0.9,0.9,0.9}{
            \makecell[{{p{\linewidth}}}]{
                \texttt{\tiny{[GM$|$GM]}}
                \texttt{Yes} \\
            }
        }
    }
    & & \\ \\

    \theutterance \stepcounter{utterance}  
    & & & \multicolumn{2}{p{0.3\linewidth}}{
        \cellcolor[rgb]{0.9,0.9,0.9}{
            \makecell[{{p{\linewidth}}}]{
                \texttt{\tiny{[GM$|$GM]}}
                \texttt{yes} \\
            }
        }
    }
    & & \\ \\

    \theutterance \stepcounter{utterance}  
    & & & \multicolumn{2}{p{0.3\linewidth}}{
        \cellcolor[rgb]{0.9,0.9,0.9}{
            \makecell[{{p{\linewidth}}}]{
                \texttt{\tiny{[GM$|$GM]}}
                \texttt{game\_result = WIN} \\
            }
        }
    }
    & & \\ \\

\end{supertabular}
}

\end{document}
